\chapter{Boils (Furuncles and Carbuncles)}

\section*{Definition}
A boil (furuncle) is a deep folliculitis that forms a tender, pus-filled nodule in the skin. A carbuncle is a cluster of interconnected furuncles. Boils are most commonly caused by Staphylococcus aureus.

\section*{When to See a Doctor}
\begin{itemize}
\item Tender, red, swollen nodule with a central pustule
\item Enlarging, painful skin lesion over several days
\item Fever, chills, or malaise accompanying a skin lesion
\item Multiple or recurrent lesions
\item Lesions in immunocompromised or diabetic people
\item Lesions on the face, near the nose or eyes, or in the groin
\item Red streaks extending from the lesion (lymphangitis)
\item Lesion that fails to improve within about two weeks
\end{itemize}

\section*{How It Develops}
Boils usually develop when bacteria, most often Staphylococcus aureus, infect a hair follicle or small break in the skin. The immune response produces inflammation and pus. A carbuncle is a connected group of boils. Recurrent lesions may reflect hidradenitis suppurativa, a local skin problem, or increased susceptibility, rather than poor hygiene alone.

\section*{Causes}
\begin{itemize}
\item Staphylococcus aureus infection (most common pathogen)
\item Methicillin-resistant Staphylococcus aureus (MRSA)
\item Streptococcal species (less common)
\item Diabetes mellitus (increased susceptibility)
\item Immunocompromised states (HIV, chemotherapy, organ transplantation)
\item Obesity (increased skin folds and friction)
\item Poor hygiene or close living conditions
\item Recurrent skin trauma or friction (shaving, tight clothing)
\item Skin conditions such as eczema or hidradenitis suppurativa
\item Contact sports with skin-to-skin transmission
\end{itemize}

\section*{Related Symptoms}
\begin{itemize}
\item Cellulitis (surrounding skin infection)
\item Abscess (larger pus collection)
\item Folliculitis (superficial hair follicle infection)
\item Impetigo (superficial skin infection)
\item Lymphangitis (red streaks from infection)
\item Fever and chills (systemic signs of infection)
\item Skin abscess
\item Hidradenitis suppurativa
\end{itemize}

\section*{Medical Evaluation}
\begin{itemize}
\item Clinical diagnosis based on physical examination
\item Wound culture and sensitivity if recurrent or suspected MRSA
\item Blood glucose testing or other tests when recurrent lesions or the clinical history suggest diabetes or another risk factor
\item Incision and drainage by a healthcare professional if large, fluctuant, or painful
\end{itemize}

\section*{Treatment Options}
\begin{itemize}
\item Incision and drainage for fluctuant abscesses (definitive treatment)
\item Warm compresses to promote spontaneous drainage
\item Antibiotics only when there is significant surrounding cellulitis, systemic illness, severe disease, impaired host defenses, or failure of drainage; the choice follows local guidance and culture when available.
\item Recurrent abscesses should be cultured and assessed for local causes. Decolonization with chlorhexidine and intranasal mupirocin is considered only in selected recurrent S. aureus cases under medical guidance.
\item Diabetes management and glycemic control to prevent recurrence
\end{itemize}

\section*{Self-Care and Home Management}
\begin{itemize}
\item Apply warm compresses to promote drainage and relieve pain
\item Keep the area clean with gentle soap and water
\item Cover the boil with a clean, dry dressing once it drains
\item Do NOT squeeze or lance the boil (risk of spreading infection)
\item Wash hands thoroughly after touching the lesion
\item Wash towels and bedding regularly; do not share towels, razors, clothing, or washcloths, and avoid swimming or contact sports until the lesion has healed.
\item Avoid sharing personal items (towels, razors, clothing)
\end{itemize}

\section*{References}
\begin{thebibliography}{99}
\bibitem{boils:ref1} Singer AJ, Talan DA. ``Management of skin abscesses in the era of methicillin-resistant Staphylococcus aureus.'' \textit{N Engl J Med}. 2014;370(11):1039-1047.
\bibitem{boils:ref2} Stevens DL, Bisno AL, et al. ``Practice guidelines for the diagnosis and management of skin and soft tissue infections.'' \textit{Clin Infect Dis}. 2014;59(2):e10-e52.
\bibitem{boils:ref3} Habif TP. \textit{Clinical Dermatology}, 6th ed. Elsevier, 2016.
\bibitem{boils:ref4} James WD, Elston DM, et al. \textit{Andrews\' Diseases of the Skin: Clinical Dermatology}, 13th ed. Elsevier, 2019.
\bibitem{boils:ref5} National Health Service. ``Boils.'' Reviewed June 2023. https://www.nhs.uk/conditions/boils/.
\bibitem{boils:ref6} Infectious Diseases Society of America. ``Practice Guidelines for the Diagnosis and Management of Skin and Soft Tissue Infections.'' 2014. https://www.idsociety.org/practice-guideline/skin-and-soft-tissue-infections/.
\end{thebibliography}
